\subsubsection{Transformation Elasticity and Capacity Utilization Estimates under a Heterogeneous Capital Framework}

The inclusion of heterogeneous capital reveals a structural paradox in the post-war US economy: while the relative cheapening of machinery drove a massive shift in the physical composition of capital, the capacity-generating power of this mechanization was severely degraded by the distributive conflicts of the 1970s. Figure \ref{fig:US_fig432_dual_theta_mu_specB} plots the implicit capacity utilization rate ($\mu_t$) alongside the asset-weighted aggregate transformation elasticity ($\theta_t^{\text{Aggregate}}$) recovered from Specification B (Model M2.44b). Unlike the static baseline elasticity of Specification A, the aggregate elasticity here is a dynamic frontier determined by the evolving mix of machinery and structures, and the macroeconomic environment mediating their installation.

As the physical accumulation of machinery and equipment accelerated relative to nonresidential structures, the aggregate transformation elasticity collapsed, severing the Harrodian link between accumulation and potential output. This structural degradation is not an artifact of index-number aggregation but the direct result of shifting marginal capacity yields. To unpack the mechanics of this collapse, Figure \ref{fig:US_fig4323_dual_axis_driving_forces} illustrates the two primary institutional driving forces behind this choice of technique: the secular decline in relative asset prices ($p_t^{\kappa} = P_t^{\text{ME}}/P_t^{\text{NRC}}$) and the episodic spikes in effective distributive pressure ($C_t = \omega_{t-1}(1-\iota_{t-1})$).

Following Stephen Marglin's pedagogical approach, we can observe the intuitive geometry of these forces. Consider a stylized benchmark during the Golden Age (e.g., the 1950s) where the structural yield of nonresidential structures is $\theta_{\text{NRC}}^{\text{str}} = 0.60$ and the equipment share is $s_{\text{ME}} = 0.30$. Under a stable distributive environment, if the marginal capacity yield of mechanization is $MP_{\kappa} = 0.40$, the aggregate elasticity is an efficient $0.72$ ($\theta^{\text{Aggregate}} = 0.60 + 0.30 \times 0.40$). However, as the post-war era progressed, machinery became structurally cheaper (the collapse of $p_t^{\kappa}$ in Figure \ref{fig:US_fig4323_dual_axis_driving_forces}), pushing the physical equipment share up toward $0.50$. 

Simultaneously, the Profit Squeeze (1968--1973) produced a sharp spike in effective distributive pressure ($C_t$). Because the interaction term $\theta_{\omega\kappa}$ is negative, this heightened wage pressure forced firms into defensive, labor-saving mechanization (Marx-biased technical change). Consequently, the marginal capacity yield of this machinery dropped significantly (e.g., down to $MP_{\kappa} = 0.20$). Thus, in the post-Fordist era, the aggregate elasticity fell to $0.70$ ($\theta^{\text{Aggregate}} = 0.60 + 0.50 \times 0.20$). Even though the physical share of machinery increased by 20 percentage points due to relative cheapening, the capacity-generating power of that machinery collapsed under the weight of defensive investment strategies induced by the profit squeeze. 

This mechanism explicitly bounds the stagnation tendency regime. Firms invested heavily in mechanization not to expand capacity, but to shed increasingly expensive labor, resulting in accumulation that was "unproductive" from the standpoint of potential output generation. The divergence between the physical volume of mechanization and its capacity yield underscores the limitations of assuming a strictly proportional conversion of capital into capacity, demonstrating why the imposition of the $\theta=1$ constraint in mainstream filtering routines consistently misreads structural shifts as cyclical deviations.

\begin{figure}[H]
    \centering
    \includegraphics[width=\textwidth]{fig/US_fig432_dual_theta_mu_specB.pdf}
    \caption{Implicit Capacity Utilization ($\mu_t$) and Asset-Weighted Transformation Elasticity ($\theta_t^{\text{Aggregate}}$)}
    \label{fig:US_fig432_dual_theta_mu_specB}
    
    \vspace{0.2cm}
    \parbox{\textwidth}{\footnotesize \textit{Note:} The figure plots the reconstructed capacity utilization rate anchored at the 1973 pinch year alongside the dynamic aggregate transformation elasticity derived from Specification B. The Harrodian knife-edge ($\theta=1.00$) and the Specification A baseline ($\theta=0.58$) are provided as reference lines.}
\end{figure}

\begin{figure}[H]
    \centering
    \includegraphics[width=\textwidth]{fig/US_fig4323_dual_axis_driving_forces.pdf}
    \caption{Effective Distributive Pressure ($C_t$) and Relative Asset Prices ($p_t^{\kappa}$)}
    \label{fig:US_fig4323_dual_axis_driving_forces}
    
    \vspace{0.2cm}
    \parbox{\textwidth}{\footnotesize \textit{Note:} The figure illustrates the institutional driving forces behind the choice of technique. The left axis tracks the effective distributive pressure to mechanize ($C_t$), while the right axis tracks the relative price of machinery and equipment to nonresidential structures ($p_t^{\kappa}$, indexed to $1973 = 1.00$).}
\end{figure}
